%Version 3.1 December 2024
% See section 11 of the User Manual for version history
%
%%%%%%%%%%%%%%%%%%%%%%%%%%%%%%%%%%%%%%%%%%%%%%%%%%%%%%%%%%%%%%%%%%%%%%
%%                                                                 %%
%% Please do not use \input{...} to include other tex files.       %%
%% Submit your LaTeX manuscript as one .tex document.              %%
%%                                                                 %%
%% All additional figures and files should be attached             %%
%% separately and not embedded in the \TeX\ document itself.       %%
%%                                                                 %%
%%%%%%%%%%%%%%%%%%%%%%%%%%%%%%%%%%%%%%%%%%%%%%%%%%%%%%%%%%%%%%%%%%%%%

%%\documentclass[referee,sn-basic]{sn-jnl}% referee option is meant for double line spacing

%%=======================================================%%
%% to print line numbers in the margin use lineno option %%
%%=======================================================%%

%%\documentclass[lineno,pdflatex,sn-basic]{sn-jnl}% Basic Springer Nature Reference Style/Chemistry Reference Style

%%=========================================================================================%%
%% the documentclass is set to pdflatex as default. You can delete it if not appropriate.  %%
%%=========================================================================================%%

%%\documentclass[sn-basic]{sn-jnl}% Basic Springer Nature Reference Style/Chemistry Reference Style

%%Note: the following reference styles support Namedate and Numbered referencing. By default the style follows the most common style. To switch between the options you can add or remove “Numbered” in the optional parenthesis. 
%%The option is available for: sn-basic.bst, sn-chicago.bst%  
 
%%\documentclass[pdflatex,sn-nature]{sn-jnl}% Style for submissions to Nature Portfolio journals
%%\documentclass[pdflatex,sn-basic]{sn-jnl}% Basic Springer Nature Reference Style/Chemistry Reference Style
\documentclass[pdflatex,sn-mathphys-num]{sn-jnl}% Math and Physical Sciences Numbered Reference Style
%%\documentclass[pdflatex,sn-mathphys-ay]{sn-jnl}% Math and Physical Sciences Author Year Reference Style
%%\documentclass[pdflatex,sn-aps]{sn-jnl}% American Physical Society (APS) Reference Style
%%\documentclass[pdflatex,sn-vancouver-num]{sn-jnl}% Vancouver Numbered Reference Style
%%\documentclass[pdflatex,sn-vancouver-ay]{sn-jnl}% Vancouver Author Year Reference Style
%%\documentclass[pdflatex,sn-apa]{sn-jnl}% APA Reference Style
%%\documentclass[pdflatex,sn-chicago]{sn-jnl}% Chicago-based Humanities Reference Style

%%%% Standard Packages
%%<additional latex packages if required can be included here>

\usepackage{graphicx}%
\usepackage{multirow}%
\usepackage{amsmath,amssymb,amsfonts}%
\usepackage{amsthm}%
\usepackage{mathrsfs}%
\usepackage[title]{appendix}%
\usepackage{xcolor}%
\usepackage{textcomp}%
\usepackage{manyfoot}%
\usepackage{booktabs}%
\usepackage{algorithm}%
\usepackage{algorithmicx}%
\usepackage{algpseudocode}%
\usepackage{listings}%
%%%%

%%%%%=============================================================================%%%%
%%%%  Remarks: This template is provided to aid authors with the preparation
%%%%  of original research articles intended for submission to journals published 
%%%%  by Springer Nature. The guidance has been prepared in partnership with 
%%%%  production teams to conform to Springer Nature technical requirements. 
%%%%  Editorial and presentation requirements differ among journal portfolios and 
%%%%  research disciplines. You may find sections in this template are irrelevant 
%%%%  to your work and are empowered to omit any such section if allowed by the 
%%%%  journal you intend to submit to. The submission guidelines and policies 
%%%%  of the journal take precedence. A detailed User Manual is available in the 
%%%%  template package for technical guidance.
%%%%%=============================================================================%%%%

%% as per the requirement new theorem styles can be included as shown below
\theoremstyle{thmstyleone}%
\newtheorem{theorem}{Theorem}%  meant for continuous numbers
%%\newtheorem{theorem}{Theorem}[section]% meant for sectionwise numbers
%% optional argument [theorem] produces theorem numbering sequence instead of independent numbers for Proposition
\newtheorem{proposition}[theorem]{Proposition}% 
%%\newtheorem{proposition}{Proposition}% to get separate numbers for theorem and proposition etc.

\theoremstyle{thmstyletwo}%
\newtheorem{example}{Example}%
\newtheorem{remark}{Remark}%

\theoremstyle{thmstylethree}%
\newtheorem{definition}{Definition}%

\raggedbottom
%%\unnumbered% uncomment this for unnumbered level heads

\begin{document}

\title[Article Title]{Article Title}

%%=============================================================%%
%% GivenName	-> \fnm{Joergen W.}
%% Particle	-> \spfx{van der} -> surname prefix
%% FamilyName	-> \sur{Ploeg}
%% Suffix	-> \sfx{IV}
%% \author*[1,2]{\fnm{Joergen W.} \spfx{van der} \sur{Ploeg} 
%%  \sfx{IV}}\email{iauthor@gmail.com}
%%=============================================================%%

\author*[1,2]{\fnm{First} \sur{Author}}\email{iauthor@gmail.com}

\author[2,3]{\fnm{Second} \sur{Author}}\email{iiauthor@gmail.com}
\equalcont{These authors contributed equally to this work.}

\author[1,2]{\fnm{Third} \sur{Author}}\email{iiiauthor@gmail.com}
\equalcont{These authors contributed equally to this work.}

\affil*[1]{\orgdiv{Department}, \orgname{Organization}, \orgaddress{\street{Street}, \city{City}, \postcode{100190}, \state{State}, \country{Country}}}

\affil[2]{\orgdiv{Department}, \orgname{Organization}, \orgaddress{\street{Street}, \city{City}, \postcode{10587}, \state{State}, \country{Country}}}

\affil[3]{\orgdiv{Department}, \orgname{Organization}, \orgaddress{\street{Street}, \city{City}, \postcode{610101}, \state{State}, \country{Country}}}




%%================================%%
%% Sample for structured abstract %%
%%================================%%

%
\abstract{\textbf{Purpose:} 
% 
%
\textbf{Methods:} 
% 
%
\textbf{Results:} 
% 
%
\textbf{Conclusion:} }




\keywords{keyword1, Keyword2, Keyword3, Keyword4}

%%\pacs[JEL Classification]{D8, H51}

%%\pacs[MSC Classification]{35A01, 65L10, 65L12, 65L20, 65L70}

\maketitle
Purpose (2–3 sentences)
Why does this study exist?

[Disease/condition] is a significant problem in [population], leading to [consequences]. While [existing approach] is the standard, current automated methods suffer from [limitations].

Methods (4–5 sentences)
What did you build or do? Highlight arch novelty

A [framework name] is proposed for [task]. [Component A] handles [function]. [Component B] enables [function]. [Strategy] is used for [purpose].

Results (2–3 sentences)
"What did you find?"

Experiments on [dataset] demonstrate [performance outcome] compared to [baseline]. The system also provides [additional output].

Conclusion (1–2 sentences)
Why does this matter?"

The proposed framework offers a [adjective] solution for [application], with potential for [broader impact].
\section{Introduction}\label{sec1}

The Introduction section, of referenced text \cite{bib1} expands on the background of the work (some overlap with the Abstract is acceptable). The introduction should not include subheadings.

Springer Nature does not impose a strict layout as standard however authors are advised to check the individual requirements for the journal they are planning to submit to as there may be journal-level preferences. When preparing your text please also be aware that some stylistic choices are not supported in full text XML (publication version), including coloured font. These will not be replicated in the typeset article if it is accepted. 

\section{Results}\label{sec2}

Sample body text. Sample body text. Sample body text. Sample body text. Sample body text. Sample body text. Sample body text. Sample body text.

\section{This is an example for first level head---section head}\label{sec3}

\subsection{This is an example for second level head---subsection head}\label{subsec2}

\subsubsection{This is an example for third level head---subsubsection head}\label{subsubsec2}

Sample body text. Sample body text. Sample body text. Sample body text. Sample body text. Sample body text. Sample body text. Sample body text. 

\section{Equations}\label{sec4}

Equations in \LaTeX\ can either be inline or on-a-line by itself (``display equations''). For
inline equations use the \verb+$...$+ commands. E.g.: The equation
$H\psi = E \psi$ is written via the command \verb+$H \psi = E \psi$+.

For display equations (with auto generated equation numbers)
one can use the equation or align environments:
\begin{equation}
\|\tilde{X}(k)\|^2 \leq\frac{\sum\limits_{i=1}^{p}\left\|\tilde{Y}_i(k)\right\|^2+\sum\limits_{j=1}^{q}\left\|\tilde{Z}_j(k)\right\|^2 }{p+q}.\label{eq1}
\end{equation}
where,
\begin{align}
D_\mu &=  \partial_\mu - ig \frac{\lambda^a}{2} A^a_\mu \nonumber \\
F^a_{\mu\nu} &= \partial_\mu A^a_\nu - \partial_\nu A^a_\mu + g f^{abc} A^b_\mu A^a_\nu \label{eq2}
\end{align}
Notice the use of \verb+\nonumber+ in the align environment at the end
of each line, except the last, so as not to produce equation numbers on
lines where no equation numbers are required. The \verb+\label{}+ command
should only be used at the last line of an align environment where
\verb+\nonumber+ is not used.
\begin{equation}
Y_\infty = \left( \frac{m}{\textrm{GeV}} \right)^{-3}
    \left[ 1 + \frac{3 \ln(m/\textrm{GeV})}{15}
    + \frac{\ln(c_2/5)}{15} \right]
\end{equation}
The class file also supports the use of \verb+\mathbb{}+, \verb+\mathscr{}+ and
\verb+\mathcal{}+ commands. As such \verb+\mathbb{R}+, \verb+\mathscr{R}+
and \verb+\mathcal{R}+ produces $\mathbb{R}$, $\mathscr{R}$ and $\mathcal{R}$
respectively (refer Subsubsection~\ref{subsubsec2}).

\section{Tables}\label{sec5}

Tables can be inserted via the normal table and tabular environment. To put
footnotes inside tables you should use \verb+\footnotetext[]{...}+ tag.
The footnote appears just below the table itself (refer Tables~\ref{tab1} and \ref{tab2}). 
For the corresponding footnotemark use \verb+\footnotemark[...]+

\begin{table}[h]
\caption{Caption text}\label{tab1}%
\begin{tabular}{@{}llll@{}}
\toprule
Column 1 & Column 2  & Column 3 & Column 4\\
\midrule
row 1    & data 1   & data 2  & data 3  \\
row 2    & data 4   & data 5\footnotemark[1]  & data 6  \\
row 3    & data 7   & data 8  & data 9\footnotemark[2]  \\
\botrule
\end{tabular}
\footnotetext{Source: This is an example of table footnote. This is an example of table footnote.}
\footnotetext[1]{Example for a first table footnote. This is an example of table footnote.}
\footnotetext[2]{Example for a second table footnote. This is an example of table footnote.}
\end{table}

\noindent
The input format for the above table is as follows:

%%=============================================%%
%% For presentation purpose, we have included  %%
%% \bigskip command. Please ignore this.       %%
%%=============================================%%
\bigskip
\begin{verbatim}
\begin{table}[<placement-specifier>]
\caption{<table-caption>}\label{<table-label>}%
\begin{tabular}{@{}llll@{}}
\toprule
Column 1 & Column 2 & Column 3 & Column 4\\
\midrule
row 1 & data 1 & data 2	 & data 3 \\
row 2 & data 4 & data 5\footnotemark[1] & data 6 \\
row 3 & data 7 & data 8	 & data 9\footnotemark[2]\\
\botrule
\end{tabular}
\footnotetext{Source: This is an example of table footnote. 
This is an example of table footnote.}
\footnotetext[1]{Example for a first table footnote.
This is an example of table footnote.}
\footnotetext[2]{Example for a second table footnote. 
This is an example of table footnote.}
\end{table}
\end{verbatim}
\bigskip
%%=============================================%%
%% For presentation purpose, we have included  %%
%% \bigskip command. Please ignore this.       %%
%%=============================================%%

\begin{table}[h]
\caption{Example of a lengthy table which is set to full textwidth}\label{tab2}
\begin{tabular*}{\textwidth}{@{\extracolsep\fill}lcccccc}
\toprule%
& \multicolumn{3}{@{}c@{}}{Element 1\footnotemark[1]} & \multicolumn{3}{@{}c@{}}{Element 2\footnotemark[2]} \\\cmidrule{2-4}\cmidrule{5-7}%
Project & Energy & $\sigma_{calc}$ & $\sigma_{expt}$ & Energy & $\sigma_{calc}$ & $\sigma_{expt}$ \\
\midrule
Element 3  & 990 A & 1168 & $1547\pm12$ & 780 A & 1166 & $1239\pm100$\\
Element 4  & 500 A & 961  & $922\pm10$  & 900 A & 1268 & $1092\pm40$\\
\botrule
\end{tabular*}
\footnotetext{Note: This is an example of table footnote. This is an example of table footnote this is an example of table footnote this is an example of~table footnote this is an example of table footnote.}
\footnotetext[1]{Example for a first table footnote.}
\footnotetext[2]{Example for a second table footnote.}
\end{table}

In case of double column layout, tables which do not fit in single column width should be set to full text width. For this, you need to use \verb+\begin{table*}+ \verb+...+ \verb+\end{table*}+ instead of \verb+\begin{table}+ \verb+...+ \verb+\end{table}+ environment. Lengthy tables which do not fit in textwidth should be set as rotated table. For this, you need to use \verb+\begin{sidewaystable}+ \verb+...+ \verb+\end{sidewaystable}+ instead of \verb+\begin{table*}+ \verb+...+ \verb+\end{table*}+ environment. This environment puts tables rotated to single column width. For tables rotated to double column width, use \verb+\begin{sidewaystable*}+ \verb+...+ \verb+\end{sidewaystable*}+.

\begin{sidewaystable}
\caption{Tables which are too long to fit, should be written using the ``sidewaystable'' environment as shown here}\label{tab3}
\begin{tabular*}{\textheight}{@{\extracolsep\fill}lcccccc}
\toprule%
& \multicolumn{3}{@{}c@{}}{Element 1\footnotemark[1]}& \multicolumn{3}{@{}c@{}}{Element\footnotemark[2]} \\\cmidrule{2-4}\cmidrule{5-7}%
Projectile & Energy	& $\sigma_{calc}$ & $\sigma_{expt}$ & Energy & $\sigma_{calc}$ & $\sigma_{expt}$ \\
\midrule
Element 3 & 990 A & 1168 & $1547\pm12$ & 780 A & 1166 & $1239\pm100$ \\
Element 4 & 500 A & 961  & $922\pm10$  & 900 A & 1268 & $1092\pm40$ \\
Element 5 & 990 A & 1168 & $1547\pm12$ & 780 A & 1166 & $1239\pm100$ \\
Element 6 & 500 A & 961  & $922\pm10$  & 900 A & 1268 & $1092\pm40$ \\
\botrule
\end{tabular*}
\footnotetext{Note: This is an example of table footnote this is an example of table footnote this is an example of table footnote this is an example of~table footnote this is an example of table footnote.}
\footnotetext[1]{This is an example of table footnote.}
\end{sidewaystable}

\section{Figures}\label{sec6}

As per the \LaTeX\ standards you need to use eps images for \LaTeX\ compilation and \verb+pdf/jpg/png+ images for \verb+PDFLaTeX+ compilation. This is one of the major difference between \LaTeX\ and \verb+PDFLaTeX+. Each image should be from a single input .eps/vector image file. Avoid using subfigures. The command for inserting images for \LaTeX\ and \verb+PDFLaTeX+ can be generalized. The package used to insert images in \verb+LaTeX/PDFLaTeX+ is the graphicx package. Figures can be inserted via the normal figure environment as shown in the below example:

%%=============================================%%
%% For presentation purpose, we have included  %%
%% \bigskip command. Please ignore this.       %%
%%=============================================%%
\bigskip
\begin{verbatim}
\begin{figure}[<placement-specifier>]
\centering
\includegraphics{<eps-file>}
\caption{<figure-caption>}\label{<figure-label>}
\end{figure}
\end{verbatim}
\bigskip
%%=============================================%%
%% For presentation purpose, we have included  %%
%% \bigskip command. Please ignore this.       %%
%%=============================================%%

\begin{figure}[h]
\centering
\includegraphics[width=0.9\textwidth]{fig.eps}
\caption{This is a widefig. This is an example of long caption this is an example of long caption  this is an example of long caption this is an example of long caption}\label{fig1}
\end{figure}

In case of double column layout, the above format puts figure captions/images to single column width. To get spanned images, we need to provide \verb+\begin{figure*}+ \verb+...+ \verb+\end{figure*}+.

For sample purpose, we have included the width of images in the optional argument of \verb+\includegraphics+ tag. Please ignore this. 

\section{Algorithms, Program codes and Listings}\label{sec7}

Packages \verb+algorithm+, \verb+algorithmicx+ and \verb+algpseudocode+ are used for setting algorithms in \LaTeX\ using the format:

%%=============================================%%
%% For presentation purpose, we have included  %%
%% \bigskip command. Please ignore this.       %%
%%=============================================%%
\bigskip
\begin{verbatim}
\begin{algorithm}
\caption{<alg-caption>}\label{<alg-label>}
\begin{algorithmic}[1]
. . .
\end{algorithmic}
\end{algorithm}
\end{verbatim}
\bigskip
%%=============================================%%
%% For presentation purpose, we have included  %%
%% \bigskip command. Please ignore this.       %%
%%=============================================%%

You may refer above listed package documentations for more details before setting \verb+algorithm+ environment. For program codes, the ``verbatim'' package is required and the command to be used is \verb+\begin{verbatim}+ \verb+...+ \verb+\end{verbatim}+. 

Similarly, for \verb+listings+, use the \verb+listings+ package. \verb+\begin{lstlisting}+ \verb+...+ \verb+\end{lstlisting}+ is used to set environments similar to \verb+verbatim+ environment. Refer to the \verb+lstlisting+ package documentation for more details.

A fast exponentiation procedure:

\lstset{texcl=true,basicstyle=\small\sf,commentstyle=\small\rm,mathescape=true,escapeinside={(*}{*)}}
\begin{lstlisting}
begin
  for $i:=1$ to $10$ step $1$ do
      expt($2,i$);  
      newline() od                (*\textrm{Comments will be set flush to the right margin}*)
where
proc expt($x,n$) $\equiv$
  $z:=1$;
  do if $n=0$ then exit fi;
     do if odd($n$) then exit fi;                 
        comment: (*\textrm{This is a comment statement;}*)
        $n:=n/2$; $x:=x*x$ od;
     { $n>0$ };
     $n:=n-1$; $z:=z*x$ od;
  print($z$). 
end
\end{lstlisting}

\begin{algorithm}
\caption{Calculate $y = x^n$}\label{algo1}
\begin{algorithmic}[1]
\Require $n \geq 0 \vee x \neq 0$
\Ensure $y = x^n$ 
\State $y \Leftarrow 1$
\If{$n < 0$}\label{algln2}
        \State $X \Leftarrow 1 / x$
        \State $N \Leftarrow -n$
\Else
        \State $X \Leftarrow x$
        \State $N \Leftarrow n$
\EndIf
\While{$N \neq 0$}
        \If{$N$ is even}
            \State $X \Leftarrow X \times X$
            \State $N \Leftarrow N / 2$
        \Else[$N$ is odd]
            \State $y \Leftarrow y \times X$
            \State $N \Leftarrow N - 1$
        \EndIf
\EndWhile
\end{algorithmic}
\end{algorithm}

%%=============================================%%
%% For presentation purpose, we have included  %%
%% \bigskip command. Please ignore this.       %%
%%=============================================%%
\bigskip
\begin{minipage}{\hsize}%
\lstset{frame=single,framexleftmargin=-1pt,framexrightmargin=-17pt,framesep=12pt,linewidth=0.98\textwidth,language=pascal}% Set your language (you can change the language for each code-block optionally)
%%% Start your code-block
\begin{lstlisting}
for i:=maxint to 0 do
begin
{ do nothing }
end;
Write('Case insensitive ');
Write('Pascal keywords.');
\end{lstlisting}
\end{minipage}

\section{Cross referencing}\label{sec8}

Environments such as figure, table, equation and align can have a label
declared via the \verb+\label{#label}+ command. For figures and table
environments use the \verb+\label{}+ command inside or just
below the \verb+\caption{}+ command. You can then use the
\verb+\ref{#label}+ command to cross-reference them. As an example, consider
the label declared for Figure~\ref{fig1} which is
\verb+\label{fig1}+. To cross-reference it, use the command 
\verb+Figure \ref{fig1}+, for which it comes up as
``Figure~\ref{fig1}''. 

To reference line numbers in an algorithm, consider the label declared for the line number 2 of Algorithm~\ref{algo1} is \verb+\label{algln2}+. To cross-reference it, use the command \verb+\ref{algln2}+ for which it comes up as line~\ref{algln2} of Algorithm~\ref{algo1}.

\subsection{Details on reference citations}\label{subsec7}

Standard \LaTeX\ permits only numerical citations. To support both numerical and author-year citations this template uses \verb+natbib+ \LaTeX\ package. For style guidance please refer to the template user manual.

Here is an example for \verb+\cite{...}+: \cite{bib1}. Another example for \verb+\citep{...}+: \citep{bib2}. For author-year citation mode, \verb+\cite{...}+ prints Jones et al. (1990) and \verb+\citep{...}+ prints (Jones et al., 1990).

All cited bib entries are printed at the end of this article: \cite{bib3}, \cite{bib4}, \cite{bib5}, \cite{bib6}, \cite{bib7}, \cite{bib8}, \cite{bib9}, \cite{bib10}, \cite{bib11}, \cite{bib12} and \cite{bib13}.

\section{Examples for theorem like environments}\label{sec10}

For theorem like environments, we require \verb+amsthm+ package. There are three types of predefined theorem styles exists---\verb+thmstyleone+, \verb+thmstyletwo+ and \verb+thmstylethree+ 

%%=============================================%%
%% For presentation purpose, we have included  %%
%% \bigskip command. Please ignore this.       %%
%%=============================================%%
\bigskip
\begin{tabular}{|l|p{19pc}|}
\hline
\verb+thmstyleone+ & Numbered, theorem head in bold font and theorem text in italic style \\\hline
\verb+thmstyletwo+ & Numbered, theorem head in roman font and theorem text in italic style \\\hline
\verb+thmstylethree+ & Numbered, theorem head in bold font and theorem text in roman style \\\hline
\end{tabular}
\bigskip
%%=============================================%%
%% For presentation purpose, we have included  %%
%% \bigskip command. Please ignore this.       %%
%%=============================================%%

For mathematics journals, theorem styles can be included as shown in the following examples:

\begin{theorem}[Theorem subhead]\label{thm1}
Example theorem text. Example theorem text. Example theorem text. Example theorem text. Example theorem text. 
Example theorem text. Example theorem text. Example theorem text. Example theorem text. Example theorem text. 
Example theorem text. 
\end{theorem}

Sample body text. Sample body text. Sample body text. Sample body text. Sample body text. Sample body text. Sample body text. Sample body text.

\begin{proposition}
Example proposition text. Example proposition text. Example proposition text. Example proposition text. Example proposition text. 
Example proposition text. Example proposition text. Example proposition text. Example proposition text. Example proposition text. 
\end{proposition}

Sample body text. Sample body text. Sample body text. Sample body text. Sample body text. Sample body text. Sample body text. Sample body text.

\begin{example}
Phasellus adipiscing semper elit. Proin fermentum massa
ac quam. Sed diam turpis, molestie vitae, placerat a, molestie nec, leo. Maecenas lacinia. Nam ipsum ligula, eleifend
at, accumsan nec, suscipit a, ipsum. Morbi blandit ligula feugiat magna. Nunc eleifend consequat lorem. 
\end{example}

Sample body text. Sample body text. Sample body text. Sample body text. Sample body text. Sample body text. Sample body text. Sample body text.

\begin{remark}
Phasellus adipiscing semper elit. Proin fermentum massa
ac quam. Sed diam turpis, molestie vitae, placerat a, molestie nec, leo. Maecenas lacinia. Nam ipsum ligula, eleifend
at, accumsan nec, suscipit a, ipsum. Morbi blandit ligula feugiat magna. Nunc eleifend consequat lorem. 
\end{remark}

Sample body text. Sample body text. Sample body text. Sample body text. Sample body text. Sample body text. Sample body text. Sample body text.

\begin{definition}[Definition sub head]
Example definition text. Example definition text. Example definition text. Example definition text. Example definition text. Example definition text. Example definition text. Example definition text. 
\end{definition}

Additionally a predefined ``proof'' environment is available: \verb+\begin{proof}+ \verb+...+ \verb+\end{proof}+. This prints a ``Proof'' head in italic font style and the ``body text'' in roman font style with an open square at the end of each proof environment. 

\begin{proof}
Example for proof text. Example for proof text. Example for proof text. Example for proof text. Example for proof text. Example for proof text. Example for proof text. Example for proof text. Example for proof text. Example for proof text. 
\end{proof}

Sample body text. Sample body text. Sample body text. Sample body text. Sample body text. Sample body text. Sample body text. Sample body text.

\begin{proof}[Proof of Theorem~{\upshape\ref{thm1}}]
Example for proof text. Example for proof text. Example for proof text. Example for proof text. Example for proof text. Example for proof text. Example for proof text. Example for proof text. Example for proof text. Example for proof text. 
\end{proof}

\noindent
For a quote environment, use \verb+\begin{quote}...\end{quote}+
\begin{quote}
Quoted text example. Aliquam porttitor quam a lacus. Praesent vel arcu ut tortor cursus volutpat. In vitae pede quis diam bibendum placerat. Fusce elementum
convallis neque. Sed dolor orci, scelerisque ac, dapibus nec, ultricies ut, mi. Duis nec dui quis leo sagittis commodo.
\end{quote}

Sample body text. Sample body text. Sample body text. Sample body text. Sample body text (refer Figure~\ref{fig1}). Sample body text. Sample body text. Sample body text (refer Table~\ref{tab3}). 

\section{Methods}\label{sec11}

This study adopts a structured end-to-end methodological pipeline for detecting sleep disorders from overnight polysomnography (PSG) recordings. The proposed framework was iteratively developed and refined through multiple experimental trials exploring different strategies for signal preprocessing, feature extraction, sequence modelling, explainability, and classification. The overall pipeline is organized into five primary modules. First, the data processing module performs channel selection, signal preprocessing, and data representation tailored for PSG time-series analysis. Second, the temporal feature extraction module learns discriminative representations from physiological signals while addressing limitations observed in prior work. Third, the sequence modelling module employs a lightweight attention-based transformer with contextual memory to capture temporal dependencies between sleep epochs and perform sleep staging. Fourth, the disorder classification module analyzes features aggregated across the full-night recording to identify the presence of sleep disorders. Finally, an explainable inference module provides interpretability by localizing important signal regions within epochs using attention-based importance mapping, estimating prediction uncertainty through confidence visualization, and utilizing entropy-based signaling to flag low-confidence epochs for potential manual review by clinicians. The complete pipeline is illustrated in Figure~\ref{fig:methodology}.

\subsection{Dataset Description}
This study utilizes two large-scale, publicly available, polysomnography (PSG) datasets: the Sleep-EDF Expanded Dataset\cite{sleepedfx, physiobank} and the Multi-Ethnic Study of Atherosclerosis (MESA) Sleep Dataset\cite{mesa, nsrr}. These datasets were selected to support the two-stage objective of this work: (i)  sleep stage classification from PSG signals, and (ii) association of sleep architecture with common sleep disorders. The Sleep-EDF Expanded dataset is used for training and validation of the sleep staging model, while the MESA dataset is employed for disorder-level analysis and for evaluating cross-dataset generalization.

\subsubsection{Sleep-EDF Expanded Dataset}

The Sleep-EDF Expanded dataset\cite{sleepedfx} was obtained from PhysioNet\cite{physiobank}, originally published in 2018 as version 1.0.0. The dataset consists of whole-night PSG recordings collected as part of two longitudinal studies: Sleep-Cassette (SC) and Sleep-Telemetry (ST). The SC subset includes recordings from healthy subjects without sleep-related medication, while the ST subset contains recordings from subjects with difficulty initiating sleep, with nights recorded under both placebo and temazepam conditions.

Across both subsets, the dataset comprises recordings from over 100 Caucasian adult subjects, spanning an age range of approximately 25 to 101 years, with each subject contributing one or more overnight recordings. In total, the expanded dataset contains  197 recorded nights, with each night lasting approximately 8 hours.

PSG data were acquired using standard laboratory-grade equipment, including EEG (Fpz–Cz and Pz–Oz electrode derivations), electrooculogram (EOG), and chin electromyogram (EMG) signals. Both EEG and EOG channels were sampled at 100 Hz. In this work, only the Fpz–Cz EEG channel was used to maintain consistency across datasets and to enable single-channel sleep staging.

Sleep stages were manually annotated by trained sleep experts following the Rechtschaffen and Kales (R\&K) standard. Each recording was segmented into 30-second epochs, with labels assigned to one of the following stages: Wake (W), N1, N2, N3, N4, REM, Movement Time, and Unknown. In accordance with current AASM conventions, N4 epochs were merged into N3, and non-physiological stages (Movement Time and Unknown) were excluded during preprocessing.

The resulting class distribution (in number of 30-second epochs) used in this study is as shown in Figure \ref{fig:stagedist}
\begin{figure}
    \centering
    \includegraphics[width=1\linewidth]{chart.pdf}
    \caption{Epoch wise Sleep Stage Distribution from Sleep-EDFx dataset}
    \label{fig:stagedist}
\end{figure}

This distribution reflects the natural imbalance present in nocturnal sleep architecture and motivates the use of deep learning–based models capable of learning better representations under severe class imbalance. 

\subsubsection{MESA Sleep Dataset}

The MESA Sleep dataset\cite{mesa} was accessed through the National Sleep Research Resource (NSRR)\cite{nsrr} and was released in 2016 as part of the Multi-Ethnic Study of Atherosclerosis. Unlike Sleep-EDF, MESA is a large-scale, population-based cohort designed to investigate the relationship between sleep characteristics and cardiometabolic health outcomes.

The dataset includes over 2,200 middle-aged and older adults (aged approximately 45–84 years) from diverse ethnic backgrounds, with each participant undergoing one full-night in-home PSG recording. PSG acquisition followed standardized protocols and included multi-channel EEG, EOG, EMG, respiratory airflow, oxygen saturation, and electrocardiogram (ECG) signals, sampled at frequencies ranging from 100 to 256 Hz depending on the modality.

Sleep stages in MESA were manually scored in 30-second epochs by certified sleep technologists according to AASM guidelines. In addition to sleep staging, MESA provides clinically adjudicated labels for several common sleep disorders, making it suitable for disorder-level analysis.

In this study, MESA was used exclusively for sleep disorder association and generalization testing, rather than for training the sleep staging model. The disorder-wise patient distribution considered in this work is summarized in the Figure \ref{fig:disorderdist}.

\begin{figure}
    \centering
    \includegraphics[width=1\linewidth]{Disorderdist.pdf}
    \caption{Class Distribution for Sleep Disorders from MESA Sleep dataset}
    \label{fig:disorderdist}
\end{figure}

This complementary use of Sleep-EDF for supervised sleep staging and MESA for disorder-level evaluation enables assessment of both model generalization mitigating learning bias and clinical relevance across heterogeneous recording conditions and populations.

\subsection{Data Preprocessing}\label{subsec:preprocessing}

Two distinct preprocessing pipelines were implemented to handle the different characteristics of the Sleep-EDF and MESA datasets. The Sleep-EDF pipeline produces epoch-level tensors for sleep staging model training, while the MESA pipeline computes subject-level clinical feature vectors for disorder classification.

\subsubsection{Sleep-EDF Signal Processing}

Raw PSG recordings in EDF format were processed using the MNE-Python library. From the available polysomnography channels, only the EEG Fpz--Cz derivation was retained for analysis. This single-channel approach was selected to ensure consistency across datasets and to reduce computational complexity while preserving the frequency content necessary for sleep stage discrimination.

Sleep stage annotations were loaded from separate hypnogram EDF files containing expert-scored labels. Stage mapping followed current AASM conventions: Wake (W), N1, N2, N3, and REM. Stage N4 was merged into N3, and non-physiological labels (Movement Time, Unknown) were excluded from training.

Continuous EEG signals were segmented into non-overlapping 30-second epochs, yielding 3,000 samples per epoch at the 100 Hz sampling rate. The 30-second epoch duration follows the standard Rechtschaffen and Kales scoring convention and enables direct comparison with clinical sleep staging practice.

Recordings were temporally trimmed to focus on sleep-containing segments. Wake epochs at recording boundaries were removed, retaining a 30-minute buffer before sleep onset and after final awakening. This reduces class imbalance by excluding extended waking periods unrelated to sleep architecture analysis.

Processed recordings were saved as PyTorch tensors with shape $[n_{\text{epochs}} \times 3000]$, where each row corresponds to one 30-second epoch. An index file records the tensor path and stage sequence for each recording.

\subsubsection{Spectral Feature Extraction}

For each 30-second epoch, complementary spectral features were computed from the EEG signal to provide frequency-domain information alongside raw waveform input.

Discrete wavelet transform features were extracted using the Daubechies-4 (db4) wavelet at decomposition level 5. For each coefficient level, four statistics were computed: energy, log-energy, Shannon entropy, and variance. The db4 wavelet was selected for its morphological similarity to sleep spindle and K-complex waveforms.

Short-time Fourier transform features were computed using a Hann window with 2-second window length and 50 percent overlap. Relative power was calculated for canonical frequency bands: Delta (0.5--4 Hz), Theta (4--8 Hz), Alpha (8--13 Hz), Beta (13--30 Hz), and Gamma (30--45 Hz). Additional spectral metrics include spectral entropy computed as $H = -\sum p_k \log p_k$ where $p_k$ is the normalized power at frequency bin $k$, and spectral edge frequency at 95 percent (SEF95) defined as the frequency below which 95 percent of total power resides.

Ratio features capture clinically relevant spectral relationships: delta-to-beta ratio, theta-to-alpha ratio, and combined slow-to-fast power ratio $((\text{delta} + \text{theta}) / (\text{alpha} + \text{beta}))$.

The complete spectral feature vector contains approximately 34 features per epoch, saved as separate tensor files aligned with raw signal epochs.

\subsubsection{MESA Clinical Feature Computation}

For the MESA dataset, a comprehensive clinical feature extraction pipeline was implemented to compute over 130 features across seven physiological domains from full PSG recordings.

EDF files were loaded via MNE with automatic channel identification using alias matching. High-frequency channels (EEG, ECG, EMG, EOG) with original sampling rates of 256 Hz were downsampled to 100 Hz using polyphase anti-alias filtering. Sleep stage annotations were parsed from NSRR-format XML files containing scored events with start time and duration fields.

Sleep architecture features include total sleep time, time in bed, sleep onset latency, wake after sleep onset, and sleep efficiency computed as $(\text{TST} / \text{TIB}) \times 100$. Stage-specific metrics include absolute duration and proportion for each stage, N3-to-N2 ratio, light-to-deep sleep ratio, REM latency, total stage transitions, transition rate per hour, and transition entropy.

Respiratory features follow AASM event detection criteria. Airflow signals were bandpass filtered (0.05--4 Hz) and the RMS envelope computed with 3-second RMS windows. Apneas were detected as 90 percent or greater amplitude reduction lasting 10 seconds or longer. Hypopneas were detected as 30 percent or greater reduction with associated desaturation. Events were classified as obstructive versus central based on respiratory effort channel amplitude during the event. Derived metrics include AHI, AI, HI, obstructive apnea index, central apnea index, REM-specific AHI, and mean and maximum apnea duration.

Oxygenation features from SpO2 analysis include mean saturation, nadir, standard deviation, T90 (time below 90 percent), T88 (time below 88 percent), ODI3 and ODI4 desaturation indices, and desaturation area-under-curve computed as the integral of deficit below 90 percent.

Cardiac and heart rate variability features were computed from ECG using R-peak detection via a Pan--Tompkins-inspired algorithm. Time-domain metrics include SDNN (standard deviation of NN intervals), RMSSD (root mean square of successive differences), and mean RR interval. Frequency-domain metrics include LF power (0.04--0.15 Hz), HF power (0.15--0.40 Hz), and LF/HF ratio via cubic spline interpolation of RR series to 4 Hz.

Periodic limb movement features were derived from leg EMG following AASM criteria: burst amplitude threshold of 8 times baseline, duration 0.5--10 seconds, inter-movement interval 5--90 seconds, and PLM series defined as 4 or more consecutive movements. Features include PLM count, PLMI (index per hour), burst rate, and IMI mean and standard deviation.

REM behavior disorder features analyze chin EMG during REM sleep. Tonic activity measures percentage of REM epochs with greater than 30 percent duration above 2 times NREM baseline. Phasic activity counts 0.1--5 second bursts above 4 times NREM baseline. Bruxism features detect 0.25--2 second bursts above 4 times background with bruxism episode index.

The final clinical feature matrix contains 136 features per subject after preprocessing steps including HRV zero-value replacement with median imputation, removal of near-zero variance features, and pruning of highly correlated duplicates ($|r| > 0.95$).

\subsection{Sleep Staging Architecture}\label{subsec:staging}

\subsection{Disorder Detection Models}\label{subsec:disorder}

\subsection{Explainability Methods}\label{subsec:explainability}

\subsection{Ethical Approval}

The MESA Sleep study was approved by institutional review boards at all participating sites. All participants provided written informed consent. This secondary analysis was conducted using de-identified data from the National Sleep Research Resource (NSRR) repository. The Sleep-EDF Expanded dataset is publicly available under open research licenses.

\section{Discussion}\label{sec12}

Discussions should be brief and focused. In some disciplines use of Discussion or `Conclusion' is interchangeable. It is not mandatory to use both. Some journals prefer a section `Results and Discussion' followed by a section `Conclusion'. Please refer to Journal-level guidance for any specific requirements. 

\section{Conclusion}\label{sec13}

Conclusions may be used to restate your hypothesis or research question, restate your major findings, explain the relevance and the added value of your work, highlight any limitations of your study, describe future directions for research and recommendations. 

In some disciplines use of Discussion or 'Conclusion' is interchangeable. It is not mandatory to use both. Please refer to Journal-level guidance for any specific requirements. 

\backmatter

\bmhead{Supplementary information}

If your article has accompanying supplementary file/s please state so here. 

Authors reporting data from electrophoretic gels and blots should supply the full unprocessed scans for key as part of their Supplementary information. This may be requested by the editorial team/s if it is missing.

Please refer to Journal-level guidance for any specific requirements.

\bmhead{Acknowledgements}

Acknowledgements are not compulsory. Where included they should be brief. Grant or contribution numbers may be acknowledged.

Please refer to Journal-level guidance for any specific requirements.

\section*{Declarations}

Some journals require declarations to be submitted in a standardised format. Please check the Instructions for Authors of the journal to which you are submitting to see if you need to complete this section. If yes, your manuscript must contain the following sections under the heading `Declarations':

\begin{itemize}
\item Funding
\item Conflict of interest/Competing interests (check journal-specific guidelines for which heading to use)
\item Ethics approval and consent to participate
\item Consent for publication
\item Data availability 
\item Materials availability
\item Code availability 
\item Author contribution
\end{itemize}

\noindent
If any of the sections are not relevant to your manuscript, please include the heading and write `Not applicable' for that section. 

%%===================================================%%
%% For presentation purpose, we have included        %%
%% \bigskip command. Please ignore this.             %%
%%===================================================%%
\bigskip
\begin{flushleft}%
Editorial Policies for:

\bigskip\noindent
Springer journals and proceedings: \url{https://www.springer.com/gp/editorial-policies}

\bigskip\noindent
Nature Portfolio journals: \url{https://www.nature.com/nature-research/editorial-policies}

\bigskip\noindent
\textit{Scientific Reports}: \url{https://www.nature.com/srep/journal-policies/editorial-policies}

\bigskip\noindent
BMC journals: \url{https://www.biomedcentral.com/getpublished/editorial-policies}
\end{flushleft}

\begin{appendices}

\section{Section title of first appendix}\label{secA1}

An appendix contains supplementary information that is not an essential part of the text itself but which may be helpful in providing a more comprehensive understanding of the research problem or it is information that is too cumbersome to be included in the body of the paper.

%%=============================================%%
%% For submissions to Nature Portfolio Journals %%
%% please use the heading ``Extended Data''.   %%
%%=============================================%%

%%=============================================================%%
%% Sample for another appendix section			       %%
%%=============================================================%%

%% \section{Example of another appendix section}\label{secA2}%
%% Appendices may be used for helpful, supporting or essential material that would otherwise 
%% clutter, break up or be distracting to the text. Appendices can consist of sections, figures, 
%% tables and equations etc.

\end{appendices}

%%===========================================================================================%%
%% If you are submitting to one of the Nature Portfolio journals, using the eJP submission   %%
%% system, please include the references within the manuscript file itself. You may do this  %%
%% by copying the reference list from your .bbl file, paste it into the main manuscript .tex %%
%% file, and delete the associated \verb+\bibliography+ commands.                            %%
%%===========================================================================================%%

\bibliography{sn-bibliography}% common bib file
%% if required, the content of .bbl file can be included here once bbl is generated
%%\input sn-article.bbl

\end{document}
